% GENERATED FILE. Edit canonical lessons, then run npm run generate:books.
% canonical-source-hash: fnv1a64:19b41e03d8d5feb4
% canonical-lessons: JA-C226-shitsumon, JA-C226-kotaeru, JA-C226-shiken, JA-C226-ukeru, JA-C226-ten

\chapter{Questions and Exams}
\label{ch:ja-questions-and-exams}
\hlchaptermodality{226}

\begin{chapteropening}
\textbf{By the end of this chapter:} \emph{I can talk about questions, replies and exams.}

Complete the last lesson of chapter 226: I can talk about questions, replies and exams.
\end{chapteropening}

\section[\texorpdfstring{\ja{しつもん} (shitsumon)}{しつもん (shitsumon)}]{\ja{しつもん} (shitsumon) — a question}

\label{lesson:JA-C226-shitsumon}



\begin{quote}
Before the new one: say the Japanese for to look up, then the Japanese for a dictionary.
\end{quote}

\subsection*{You'll want to know: \ja{しつもん}}
\textbf{\ja{しつもん}} — \emph{shitsumon} — \textquotedblleft{}a question\textquotedblright{}.

\begin{quote}
\textbf{\ja{がくせいは} \ja{せんせいに} \ja{しつもんを} \ja{します。}} — \emph{gakusei wa sensei ni shitsumon o shimasu} — \textquotedblleft{}The students ask the teacher questions.\textquotedblright{}
\end{quote}

\begin{cousinweb}
\textbf{\ja{しつもん}} (\emph{shitsumon}) is Sino-Japanese (kango): two Chinese characters read with their on'yomi, readings borrowed from Chinese. They are \emph{shitsu}, \textquotedblleft{}substance\textquotedblright{} (Middle Chinese \emph{tsyit}) and \emph{mon}, \textquotedblleft{}ask\textquotedblright{} (\emph{mjunH}) (KANJIDIC2; Middle Chinese in Baxter's transcription).
\end{cousinweb}

\begin{grammarlens}[title={What you've built}]
One more word for everyday life at A2.
\end{grammarlens}

\subsection*{Your turn}
\begin{itemize}
\raggedright
  \item \emph{Say it:} \emph{shitsumon}
  \item \emph{Say it:} \emph{shitsumon}, once more
  \item \emph{Say it:} \emph{gakusei wa sensei ni shitsumon o shimasu} — The students ask the teacher questions.
  \item \emph{Recall:} say the Japanese for to look up, then the Japanese for a dictionary, then say \emph{shitsumon} again
\end{itemize}

\begin{tcolorbox}[breakable,colback=teal!4,colframe=teal!35!black,title={Before you move on}]
What does \ja{しつもん} mean? (\textquotedblleft{}A question\textquotedblright{}.) Say it once more.
\end{tcolorbox}
\section[\texorpdfstring{\ja{こたえる} (kotaeru)}{こたえる (kotaeru)}]{\ja{こたえる} (kotaeru) — to reply; to answer}

\label{lesson:JA-C226-kotaeru}



\begin{quote}
Before the new one: say the Japanese for a dictionary, then the Japanese for a question.
\end{quote}

\subsection*{You'll want to know: \ja{こたえる}}
\textbf{\ja{こたえる}} — \emph{kotaeru} — \textquotedblleft{}to reply; to answer\textquotedblright{}.

\begin{quote}
\textbf{\ja{せんせいは} \ja{しつもんに} \ja{ていねいに} \ja{こたえます。}} — \emph{sensei wa shitsumon ni teinei ni kotaemasu} — \textquotedblleft{}The teacher replies to questions politely.\textquotedblright{}
\end{quote}

Polite \textbf{\ja{こたえます}}, \emph{kotaemasu}; an ichidan verb. The noun \textbf{\ja{こたえ}}, \emph{kotae}, \textquotedblleft{}an answer\textquotedblright{}, comes from it.

\begin{cousinweb}
\textbf{\ja{こたえる}} (\emph{kotaeru}) is a native Japanese word (wago), inherited rather than borrowed from Chinese; its kanji was attached later as a kun'yomi (JMdict, KANJIDIC2).
\end{cousinweb}

\begin{grammarlens}[title={What you've built}]
One more word for everyday life at A2.
\end{grammarlens}

\subsection*{Your turn}
\begin{itemize}
\raggedright
  \item \emph{Say it:} \emph{kotaeru}
  \item \emph{Say it:} \emph{kotaeru}, once more
  \item \emph{Say it:} \emph{sensei wa shitsumon ni teinei ni kotaemasu} — The teacher replies to questions politely.
  \item \emph{Recall:} say the Japanese for a dictionary, then the Japanese for a question, then say \emph{kotaeru} again
\end{itemize}

\begin{tcolorbox}[breakable,colback=teal!4,colframe=teal!35!black,title={Before you move on}]
What does \ja{こたえる} mean? (\textquotedblleft{}To reply; to answer\textquotedblright{}.) Say it once more.
\end{tcolorbox}
\section[\texorpdfstring{\ja{しけん} (shiken)}{しけん (shiken)}]{\ja{しけん} (shiken) — an exam; a test}

\label{lesson:JA-C226-shiken}



\begin{quote}
Before the new one: say the Japanese for a question, then the Japanese for to reply.
\end{quote}

\subsection*{You'll want to know: \ja{しけん}}
\textbf{\ja{しけん}} — \emph{shiken} — \textquotedblleft{}an exam; a test\textquotedblright{}.

\begin{quote}
\textbf{\ja{がっこうは} \ja{まいしゅう} \ja{しけんが} \ja{あります。}} — \emph{gakkō wa maishū shiken ga arimasu} — \textquotedblleft{}The school has an exam every week.\textquotedblright{}
\end{quote}

\begin{cousinweb}
\textbf{\ja{しけん}} (\emph{shiken}) is Sino-Japanese (kango): two Chinese characters read with their on'yomi, readings borrowed from Chinese. They are \emph{shi}, \textquotedblleft{}try\textquotedblright{} (Middle Chinese \emph{syiH}) and \emph{ken}, \textquotedblleft{}verify\textquotedblright{} (\emph{ngjemH}) (KANJIDIC2; Middle Chinese in Baxter's transcription).
\end{cousinweb}

\begin{grammarlens}[title={What you've built}]
One more word for everyday life at A2.
\end{grammarlens}

\subsection*{Your turn}
\begin{itemize}
\raggedright
  \item \emph{Say it:} \emph{shiken}
  \item \emph{Say it:} \emph{shiken}, once more
  \item \emph{Say it:} \emph{gakkō wa maishū shiken ga arimasu} — The school has an exam every week.
  \item \emph{Recall:} say the Japanese for a question, then the Japanese for to reply, then say \emph{shiken} again
\end{itemize}

\begin{tcolorbox}[breakable,colback=teal!4,colframe=teal!35!black,title={Before you move on}]
What does \ja{しけん} mean? (\textquotedblleft{}An exam; a test\textquotedblright{}.) Say it once more.
\end{tcolorbox}
\section[\texorpdfstring{\ja{うける} (ukeru)}{うける (ukeru)}]{\ja{うける} (ukeru) — to take an exam; to receive}

\label{lesson:JA-C226-ukeru}



\begin{quote}
Before the new one: say the Japanese for to reply, then the Japanese for an exam.
\end{quote}

\subsection*{You'll want to know: \ja{うける}}
\textbf{\ja{うける}} — \emph{ukeru} — \textquotedblleft{}to take an exam; to receive\textquotedblright{}.

\begin{quote}
\textbf{\ja{あには} \ja{らいねん} \ja{しけんを} \ja{うけます。}} — \emph{ani wa rainen shiken o ukemasu} — \textquotedblleft{}My elder brother takes an exam next year.\textquotedblright{}
\end{quote}

Polite \textbf{\ja{うけます}}, \emph{ukemasu}; an ichidan verb.

\begin{cousinweb}
\textbf{\ja{うける}} (\emph{ukeru}) is a native Japanese word (wago), inherited rather than borrowed from Chinese; its kanji was attached later as a kun'yomi (JMdict, KANJIDIC2).
\end{cousinweb}

\begin{grammarlens}[title={What you've built}]
One more word for everyday life at A2.
\end{grammarlens}

\subsection*{Your turn}
\begin{itemize}
\raggedright
  \item \emph{Say it:} \emph{ukeru}
  \item \emph{Say it:} \emph{ukeru}, once more
  \item \emph{Say it:} \emph{ani wa rainen shiken o ukemasu} — My elder brother takes an exam next year.
  \item \emph{Recall:} say the Japanese for to reply, then the Japanese for an exam, then say \emph{ukeru} again
\end{itemize}

\begin{tcolorbox}[breakable,colback=teal!4,colframe=teal!35!black,title={Before you move on}]
What does \ja{うける} mean? (\textquotedblleft{}To take an exam; to receive\textquotedblright{}.) Say it once more.
\end{tcolorbox}
\section[\texorpdfstring{\ja{てん} (ten)}{てん (ten)}]{\ja{てん} (ten) — a mark; a score; a point}

\label{lesson:JA-C226-ten}



\begin{quote}
Before the new one: say the Japanese for an exam, then the Japanese for to take an exam.
\end{quote}

\subsection*{You'll want to know: \ja{てん}}
\textbf{\ja{てん}} — \emph{ten} — \textquotedblleft{}a mark; a score; a point\textquotedblright{}.

\begin{quote}
\textbf{\ja{あにの} \ja{しけんの} \ja{てんは} \ja{いいです。}} — \emph{ani no shiken no ten wa ii desu} — \textquotedblleft{}My elder brother's exam mark is good.\textquotedblright{}
\end{quote}

\begin{cousinweb}
\textbf{\ja{てん}} (\emph{ten}) is Sino-Japanese (kango): one Chinese character read with its on'yomi, a reading borrowed from Chinese: \emph{ten}, \textquotedblleft{}dot, mark\textquotedblright{} (Middle Chinese \emph{temX}) (KANJIDIC2; Middle Chinese in Baxter's transcription).
\end{cousinweb}

\begin{grammarlens}[title={What you've built}]
One more word for everyday life at A2.
\end{grammarlens}

\subsection*{Your turn}
\begin{itemize}
\raggedright
  \item \emph{Say it:} \emph{ten}
  \item \emph{Say it:} \emph{ten}, once more
  \item \emph{Say it:} \emph{ani no shiken no ten wa ii desu} — My elder brother's exam mark is good.
  \item \emph{Recall:} say the Japanese for an exam, then the Japanese for to take an exam, then say \emph{ten} again
\end{itemize}

\begin{tcolorbox}[breakable,colback=teal!4,colframe=teal!35!black,title={Before you move on}]
What does \ja{てん} mean? (\textquotedblleft{}A mark; a score; a point\textquotedblright{}.) Say it once more.
\end{tcolorbox}
